\documentclass{bmstu}

\begin{document}

\makereporttitle
{Информатика и системы управления (ИУ)}
{Программное обеспечение ЭВМ и информационные технологии (ИУ7)}
{\textbf{3}}
{Проектирование экспертных систем}
{Обратный поиск в глубину от цели в графах И/ИЛИ}
{}
{ИУ7-32М}
{К.Э. Ковалец}
{З.Н. Русакова}


\setcounter{page}{2}
\renewcommand{\contentsname}{Содержание} 
\tableofcontents


\chapter{Основная часть}

\section{Описание метода}

Обратный поиск в глубину --- это метод, который используется для нахождения пути от заданной цели к начальным
вершинам в графе И/ИЛИ. В процессе работы используются следующие списки:
\begin{itemize}
    \item стек открытых вершин;
    \item список закрытых вершин;
    \item список запрещенных вершин;
    \item список открытых правил;
    \item список закрытых правил;
    \item список запрещенных правил.
\end{itemize}

\section{Постановка задачи}

Необходимо разработать класс метода поиска в глубину, полями которого являются указанные списки, целевая вершина, а также два флага: \texttt{solution\_flg} --- есть решение и \texttt{no\_solution\_flg} --- нет решения, имеющие значения по умолчанию false. В конструктор класса передается список базы правил. Целевая вершина и массив заданных входных вершин задаются при запуске метода поиска. Входные вершины записываются в закрытые вершины, целевая вершина помещается в голову стека (список открытых вершин). Все вершины помечаются флагами, равными нулю.

Разрабатываются четыре метода:
\begin{itemize}
    \item <<поиск потомков>> --- определяет первое правило, раскрывающее текущую подцель, и возвращает, было ли найдено правило;
    \item <<поиск>> (основной метод) --- вызывает все методы алгоритма, работает в цикле while, пока флаги равны false;
    \item <<возврат>> (back tracking) --- отсекает запрещённые вершины и правила;
    \item <<разметка>> --- формирует список закрытых правил, составляющих дерево решения, и список закрытых (доказанных) вершин.
\end{itemize}

\section{Алгоритм поиска потомков}

Текущая подцель выбирается из головы стека открытых вершин. В цикле по базе правил определяется первое правило, выходная вершина которого совпадает с подцелью (раскрывает эту вершину). Номер правила записывается в голову стека открытых правил. Затем определяются новые подцели и записываются в стек открытых вершин. Для этого необходимо определить, какие вершины выбранного правила из входных не входят в закрытые. Эти вершины добавляются в голову стека открытых вершин. Если все входные вершины правила попали в закрытые, то новых подцелей нет и вызывается процедура разметки.

\section{Алгоритм возврата}

Если потомков не найдено, функция возвращает ноль. Это означает, что текущая подцель должна быть объявлена как запрещённая вершина, а само правило --- как запрещённое. Необходимо удалить текущую подцель из стека открытых вершин (из головы), объявить правило из головы стека запрещённым, удалить его и перенести в запрещённые. Также необходимо удалить из стека все вершины, которые являлись входными вершинами найденного запрещённого правила. После этого снова вызывается метод поиска потомков для нахождения альтернативного пути. Если альтернативы нет, то снова происходит возврат. Если в стеке открытых вершин осталась только одна целевая после ненахождения нужного правила, то решение не найдено (\texttt{no\_solution\_flg} = true).

\section{Алгоритм разметки}

Этот алгоритм используется для построения дерева решения и закрытия вершин и правил. Он работает путём
последовательного подъема снизу-вверх по дереву.

Основные шаги алгоритма:
\begin{itemize}
    \item извлечение правила из списка открытых правил и добавление его в список закрытых правил;
    \item извлечение вершины из списка открытых вершин и добавление ее в список закрытых вершин;
    \item проверка, является ли текущая вершина целевым узлом. Если да, то найдено решение (\texttt{solution\_flg} = true);
    \item если текущая вершина не является целевой, то необходимо вернуться к последнему открытому узлу и
    последнему открытому правилу и проверить, соответствует ли текущее правило текущему узлу. Если нет,
    то алгоритм завершается.
\end{itemize}

\section{Результаты работы программы}

Следующими цветами обозначены отдельные группы правил и вершин:
\begin{itemize}
    \item зеленым --- целевая вершина графа;
    \item синим --- входные вершины графа;
    \item фиолетовым --- закрытые вершина и правила графа, не включая входные вершины;
    \item оранжевым --- запрещенные вершины и правила графа.
\end{itemize}

Результаты работы программы приведены на рисунках \ref{img:initial_graph.gv.pdf}--\ref{img:final_graph.gv.pdf}.

\clearpage
\imgs{initial_graph.gv.pdf}{h!}{0.55}{Входной граф И/ИЛИ}

\imgs{final_graph.gv.pdf}{h!}{0.55}{Итоговый граф И/ИЛИ с построенным деревом решений}


\chapter{Текст программы}

В листингах \ref{lst:main}--\ref{lst:search} представлен код программы.

\mylisting[python]{main.py}{firstline=1,lastline=38}{Основной модуль программы}{main}{}

\mylisting[python]{data.py}{firstline=1,lastline=39}{Модуль с данными}{data}{}

\mylisting[python]{graph.py}{firstline=1,lastline=92}{Модуль для построения графа}{graph}{}

\mylisting[python]{items.py}{firstline=1,lastline=35}{Модуль с описанием правил и вершин}{items}{}

\mylisting[python]{stack.py}{firstline=1,lastline=54}{Модуль с реализацией стека}{stack}{}

\mylisting[python]{search.py}{firstline=1,lastline=160}{Модуль с реализацией обратного поиска в глубину от цели}{search}{}

\end{document}
